\subsection{Unit Vectors and Orthogonality}

\begin{definition} \label{def:unit-vector-and-normalization} \hlt{Unit Vector and Normalization}\\
A vector $\mathbf{u}\in\mathcal{V}$ is a \hlt{unit vector} if $\norm{\mathbf{u}}=1$. If $\mathbf{v}\neq\mathbf{0}$, then $\frac{\mathbf{v}}{\norm{\mathbf{v}}}$ is called the \hlt{normalization} of $\mathbf{v}$.
\end{definition}

\begin{proposition} \label{prop:normalization-gives-unit-vector} \hlt{Normalization Gives a Unit Vector}\\
If $\mathbf{v}\in\mathcal{V}$ and $\mathbf{v}\neq\mathbf{0}$, then $\frac{\mathbf{v}}{\norm{\mathbf{v}}}$ is a unit vector in the same direction as $\mathbf{v}$.
\begin{proof}
By Proposition \ref{prop:derived-norm-basic-properties}, $\norm{\mathbf{v}}>0$ and $\norm{\frac{\mathbf{v}}{\norm{\mathbf{v}}}} =\frac{1}{\norm{\mathbf{v}}}\norm{\mathbf{v}}=1$.\\
The vector $\frac{\mathbf{v}}{\norm{\mathbf{v}}}$ is a positive scalar multiple of $\mathbf{v}$, so it lies in the same one-dimensional direction.
\hlt{Structural Commentary}: Normalization changes scale but preserves direction.\\
\hlt{Key Insight}: Divide by the length to make the length equal to one.
\end{proof}
\end{proposition}

\begin{definition} \label{def:orthogonal-vectors} \hlt{Orthogonal Vectors}\\
Vectors $\mathbf{u},\mathbf{v}\in\mathcal{V}$ are \hlt{orthogonal}, written $\mathbf{u}\perp\mathbf{v}$, if $\innerprod{\mathbf{u},\mathbf{v}}=0$.
\end{definition}

\begin{proposition} \label{prop:orthogonality-and-zero} \hlt{Orthogonality and the Zero Vector}\\
Let $\mathcal{V}$ be an inner-product space.
\begin{enumerate}[label=\roman*.]
\setlength{\itemsep}{0pt}
\item $\mathbf{u}\perp\mathbf{v}$ if and only if $\mathbf{v}\perp\mathbf{u}$.
\item $\mathbf{0}$ is orthogonal to every vector in $\mathcal{V}$.
\item If $\mathbf{v}$ is orthogonal to itself, then $\mathbf{v}=\mathbf{0}$.
\end{enumerate}
\begin{proof}
\begin{enumerate}[label=(\roman*)]
\setlength{\itemsep}{0pt}
\item If $\innerprod{\mathbf{u},\mathbf{v}}=0$, then conjugate symmetry gives $\innerprod{\mathbf{v},\mathbf{u}}=\overline{0}=0$, and the converse is identical.
\item This follows from Proposition \ref{prop:basic-inner-product-identities}, which gives $\innerprod{\mathbf{0},\mathbf{u}}=0$ and $\innerprod{\mathbf{u},\mathbf{0}}=0$.
\item If $\mathbf{v}\perp\mathbf{v}$, then $\innerprod{\mathbf{v},\mathbf{v}}=0$, hence $\mathbf{v}=\mathbf{0}$ by definiteness.
\end{enumerate}
\hlt{Structural Commentary}: Conjugate symmetry makes orthogonality symmetric, and definiteness isolates $\mathbf{0}$.\\
\hlt{Key Insight}: Conjugate symmetry preserves the equation ``equals zero.''
\end{proof}
\end{proposition}

\begin{proposition} \label{prop:mutually-orthogonal-nonzero-vectors-independent} \hlt{Mutually Orthogonal Nonzero Vectors Are Independent}\\
Mutually orthogonal nonzero vectors $\mathbf{v}_1,\ldots,\mathbf{v}_m$ in an inner-product space are linearly independent.
\begin{proof}
If $\sum\limits_{j=1}^m c_j\mathbf{v}_j=\mathbf{0}$, then for each $k$, $0=\innerprod{\sum\limits_{j=1}^m c_j\mathbf{v}_j,\mathbf{v}_k} =c_k\innerprod{\mathbf{v}_k,\mathbf{v}_k}$.\\
Since $\innerprod{\mathbf{v}_k,\mathbf{v}_k}>0$, $c_k=0$. Thus all coefficients vanish.\\
\hlt{Structural Commentary}: Orthogonality isolates one coefficient at a time.\\
\hlt{Key Insight}: Pair a dependence relation with one vector from the list; every cross term disappears.
\end{proof}
\end{proposition}
